\section{3D Reconstruction}
\label{sec:reconstruction}
% ----- Camera calibration --------------------------------------
\begin{frame}{3D Reconstruction: Camera Calibration}
\begin{itemize}
  \item \textbf{World frame:} a standard FIBA court model (true court dimensions in mm; origin at court centre).
  \item \textbf{Intrinsics:} known a priori for each camera.
  \item \textbf{Extrinsics:} estimated by solving Perspective-n-Point (PnP).
  \item PnP maps manually annotated court landmarks to their 3D world coordinates.
\end{itemize}
\end{frame}

% ----- Triangulation (video) -----------------------------------
\begin{frame}{3D Reconstruction: Triangulation}
\begin{columns}[T]
  \begin{column}{0.55\textwidth}
    Players and the ball are treated differently:
    \begin{itemize}
      \item \textbf{Players:} a planar homography maps the foot pixel onto court-floor coordinates.
      \begin{itemize}
        \item Floor constraint allows single-view localisation under occlusion.
        \item Estimates fused across the three cameras after discarding outliers.
      \end{itemize}
      \item \textbf{Ball:} airborne, so triangulated from at least $2$ views.
      \begin{itemize}
        \item DLT stacks the per-camera ray constraints; SVD solves the system.
        \item High-error views discarded and the point re-triangulated.
      \end{itemize}
    \end{itemize}
  \end{column}
  \begin{column}{0.42\textwidth}
    \videobox
  \end{column}
\end{columns}
\end{frame}

% ----- Trajectory smoothing ------------------------------------
\begin{frame}{3D Reconstruction: Trajectory Smoothing}
Remove residual noise from detection jitter and calibration error.
\begin{itemize}
  \item Forward Kalman filter, then a Rauch--Tung--Striebel backward pass.
  \item Past and future observations correct each state.
  \item Reject implausible spatial jumps as outliers.
  \item Split a trajectory at long gaps, rather than interpolate across frames with no detections.
\end{itemize}
\end{frame}
